\documentclass[11pt]{amsart}
\usepackage[margin=1.15in]{geometry}
\usepackage{amssymb,amsmath,mathtools}
\usepackage{booktabs}
\usepackage{array}
\usepackage{xcolor}
\usepackage{url}

\usepackage[colorlinks=true,linkcolor=blue!50!black,citecolor=blue!50!black,urlcolor=blue!50!black]{hyperref}
\newcommand{\doi}[1]{doi:\,\href{https://doi.org/#1}{#1}}

% ------------------------------------------------------------------ macros
\newcommand{\Lmax}{1.4035}                                % right end of the certified range (release 1.1.0; goal 1.5)
\newcommand{\ev}{\mathrm{ev}}
\newcommand{\od}{\mathrm{odd}}
\newcommand{\R}{\mathbb{R}}
\newcommand{\re}{\operatorname{Re}}
\newcommand{\pole}{\mathrm{pole}}
\newcommand{\Tsh}{T^{\sharp}}
\newcommand{\eps}{\varepsilon}
\newcommand{\file}[1]{{\small\path{#1}}}
\newcommand{\e}[1]{\cdot 10^{#1}}

\newtheorem{theorem}{Theorem}[section]
\newtheorem{proposition}[theorem]{Proposition}
\newtheorem{lemma}[theorem]{Lemma}
\newtheorem{corollary}[theorem]{Corollary}
\theoremstyle{definition}
\newtheorem{definition}[theorem]{Definition}
\theoremstyle{remark}
\newtheorem{remark}[theorem]{Remark}

\numberwithin{equation}{section}

\begin{document}

\title[Simple, even ground state of the localized Weil form]{A certified simple, even ground state of the localized Weil quadratic form for all windows $L\in(0,\Lmax]$}

\author{Johann Alexander Botha}
\address{GenusCore}
\email{info@genuscore.com}

\date{Version 1.1.0, 31 October 2026}

\subjclass[2020]{11M26 (primary); 65G20, 47A75, 65F15 (secondary)}
\keywords{Weil explicit formula, Weil positivity, quadratic form, ground state, computer-assisted proof, ball arithmetic}

\begin{abstract}
Let $Q_L$ be Weil's quadratic form for the Riemann zeta function, restricted to real test functions supported in $[-L,L]$ on the logarithmic scale; it contains the pole terms, the archimedean term and every prime power $n$ with $\log n<2L$. Connes and van Suijlekom proved that if the lowest eigenvalue of the associated self-adjoint operator is simple with an even eigenfunction~$\xi$, then $\hat\xi$ has only real zeros; for the Weil form this ``simple-even'' property is one of the missing steps in the strategy of Connes, Consani and Moscovici. We give a computer-assisted proof, in Arb ball arithmetic, that the ground state of $Q_L$ is simple and even for every $L\in(0,\Lmax]$. The proof combines min-max monotonicity in~$L$ with three pieces: a Lipschitz argument on $(0,0.05]$, a covering of $[0.05,1.19]$ by $31$ certified windows based on Zhu's frequency-cutoff reduction, and a covering of $(1.19,\Lmax]$ by $8$ windows in which second-eigenvalue lower bounds are obtained from Elston's window-compressed prime-norm certifier combined with a rank-one deflation, and upper bounds from Legendre--Galerkin Rayleigh quotients with rigorous quadrature bounds. On $(0,0.8]$ the statement had been claimed before; the range $(0.8,\Lmax]$ is, to our knowledge, new. As by-products we give rigorous upper bounds for the lowest eigenvalue in both parity sectors up to $L=2.25$.
The extension of the same computation to $L=1.5$ is in progress.
\end{abstract}

\maketitle

\setcounter{tocdepth}{1}
\tableofcontents

% =====================================================================
\section{Introduction}\label{sec:intro}

\subsection{The localized Weil form}
For a real function $f\in L^2(\R)$ with support in $[-L,L]$ put $\tilde f(x)=f(-x)$, $g=f*\tilde f$ and $F(z)=\hat f(z)=\int f(x)e^{izx}\,dx$. Weil's explicit-formula functional~\cite{Weil}, evaluated at~$g$, gives the quadratic form
\begin{multline}\label{eq:weil}
Q_L(f)=W(g),\qquad
W(g)=\hat g(\tfrac i2)+\hat g(-\tfrac i2)-g(0)\log\pi+\frac1{2\pi}\int_\R\hat g(r)\,\re\psi\bigl(\tfrac14+\tfrac{ir}{2}\bigr)\,dr\\
-\sum_{n\ge2}\frac{\Lambda(n)}{\sqrt n}\bigl[g(\log n)+g(-\log n)\bigr],
\end{multline}
where $\psi=\Gamma'/\Gamma$ and $\Lambda$ is the von Mangoldt function. Weil's criterion states that the Riemann hypothesis (RH) is equivalent to $Q_L(f)\ge0$ for all such $f$ and all $L>0$~\cite{Weil,Bombieri}. Since $g$ is supported in $[-2L,2L]$, only prime powers with $\log n<2L$ contribute; the prime power $n$ enters at $L=\tfrac12\log n$ (Table~\ref{tab:primes}).

\begin{table}[ht]
\centering\small
\caption{Prime powers $n$ and the window $L=\tfrac12\log n$ at which they enter the form, for $L\le1.5$.}\label{tab:primes}
\begin{tabular}{l*{6}{c}}
\toprule
$n$ & 2 & 3 & 4 & 5 & 7 & 8\\
$\tfrac12\log n$ & 0.3466 & 0.5493 & 0.6931 & 0.8047 & 0.9730 & 1.0397\\
\midrule
$n$ & 9 & 11 & 13 & 16 & 17 & 19\\
$\tfrac12\log n$ & 1.0986 & 1.1989 & 1.2825 & 1.3863 & 1.4166 & 1.4722\\
\bottomrule
\end{tabular}
\end{table}

\subsection{The simple-even hypothesis}
Connes and van Suijlekom~\cite[Thm.~6.1]{CvS} prove: if a quadratic form with Schwartz kernel $\tilde D(x-y)$, $D$ a real distribution, defines a lower-bounded self-adjoint operator whose lowest spectral value is a \emph{simple, isolated} eigenvalue with \emph{even} eigenfunction $\xi$, then all zeros of the entire function $\hat\xi$ are real. Connes, Consani and Moscovici~\cite{CCM25} propose a strategy towards RH based on a spectral realization of the zeros through the eigenfunctions $\xi_\lambda$ of the Weil form on growing windows; in \cite[\S8]{CCM25} they list two ``missing steps'': (i) the smallest eigenvalue of the Weil quadratic form $QW_\lambda$ is simple with even eigenvector, and (ii) the zeros of $\hat\xi_\lambda$ converge to the nontrivial zeros of $\zeta$. The same two steps are listed in~\cite[\S6.6]{Connes26}. The present paper addresses only step~(i), and only on the bounded range $L\in(0,\Lmax]$; it says nothing about step~(ii) or about large~$L$, and does not prove positivity beyond what the certified lower bounds imply on this range. In the conventions of~\cite{CvS} the operator acts on $L^2([-\tfrac L2,\tfrac L2])$, so their $L$ is twice ours; in~\cite{CCM25} the window is $[\lambda^{-1},\lambda]$ on the multiplicative scale, i.e.\ $L=\log\lambda$. Our form agrees with theirs only up to a positive affine rescaling, see Remark~\ref{rem:affine} and \S\ref{sec:limitations}.

\subsection{Prior work}\label{sec:prior}
We list what we know to have been proved or claimed; status labels follow the sources. Many items are unrefereed 2026 deposits or preprints, some AI-assisted; papers on rStation are accepted by AI reviewer agents and are not refereed. Our literature search (arXiv, Zenodo, GitHub and rStation) was last refreshed on 31 October 2026.
\begin{itemize}
\item \emph{Small windows.} Suzuki~\cite[Thm.~1.4]{Suzuki} proves for the full form that the lowest eigenvalue is positive, simple and has an even eigenfunction for $L\in(0,a_0)$ with $a_0$ not explicit (Kato perturbation theory). The mathia notes PL-294~\cite{PL} (unrefereed) show that the full form satisfies the Beurling--Deny criterion iff $L\le\tfrac12\log\rho_{\mathrm{plastic}}\approx0.1406$, which gives a simple ground state of fixed sign. Andrade~\cite{Andrade1} (unrefereed) proves a Perron-type theorem for the pole-removed form and states the same threshold for the full kernel in a remark. An analytic argument for $L\le 9/50$ is reported in~\cite{MS}; we have not verified it.
\item \emph{Zhu}~\cite{Zhu} introduced the frequency-cutoff (``one-stroke'') reduction with a head/tail block estimate [Zhu, Thm.~1.1] on which all our window certificates rest, certified simple-even at the single window $L=0.8$ [Zhu, Thm.~6.2], and certified variational upper bounds on the lowest eigenvalue for $L\in[0.5,2]$ [Zhu, Table~3]. Zhu [\S7] retracts an earlier positivity claim at $L=1.19$.
\item \emph{Kramarenko-Byrd}~\cite{KB} (unrefereed) claims a rigorous proof of simple-even for all $0<L\le4/5$ in our half-width convention. We have not checked it. If it is correct, our result on $(0,0.8]$ is an independent confirmation.
\item \emph{Perturbative simplicity} on a neighbourhood of $L=0.8$ of non-explicit radius: PL-287~\cite{PL}.
\item \emph{Positivity and per-sector lower bounds} (no simplicity or parity statement): jose-paul-b~\cite{JPB} for $L\le\tfrac12\log7$; Liu~\cite{Liu} at $L=17/16$ (block-Schur bounded comparison, $Q\ge2^{-49162}$); Mori~\cite{Mori} at $L=\tfrac12\log15$; Senke~\cite{Senke} at $L=77/64$; Elder~\cite{Elder} at $L=0.8$; and Elston~\cite{Elston}, whose certified lowest-eigenvalue lower bounds at $L=1.19$, $1.3$ and $1.495$ (e.g.\ $\lambda\ge3.501\e{-96}$ at $L=1.495$) use the certifier we build on in Piece~III. On rStation, lihongmin~\cite[Thm.~1.1]{rh45} certifies strict positivity of both parity sectors for every support radius $0<a\le\tfrac12\log11$, and bramble\_trace~\cite[Thm.~1.1]{rh60} for every $0<a\le2$, both by block Schur elimination with the complete archimedean response. Their form~\cite[(2.1)]{rh45}, \cite[(2.1)]{rh60} coincides term by term with~\eqref{eq:weil}, with $a=L$. Neither paper identifies the ground state or compares the two sectors (\cite[\S1]{rh45}: ``our lower certificate neither identifies the ground state nor estimates the sharp minimum''), and \cite{rh60} gives no explicit eigenvalue lower bound.
\item \emph{Upper bounds for large $L$.} Without hypotheses on the zeros, Mori~\cite[Thm.~K(a)]{MoriV} proves $\lambda_{\min}\le3.50\times10^{24}\mu^{9/2}\log\mu\,e^{-4\pi\mu}$ for $50\le\mu=e^{2L}\le10^4$, in his normalisation, and bramble\_trace~\cite[Thm.~2.1]{rh25} proves $\lambda_1^\ev(L)\le10^{16}\exp(11L-2\pi e^{2L})$ for $L\ge1$ in the same form as ours. Under RH, span\_of\_dilates~\cite[Thm.~1]{rh81} proves $-\ln\lambda^*(L)\ge4\pi e^{2L}-O(L^2)$ in both sectors, in Zhu's normalisation. The decay scale $e^{-4\pi e^{2L}}$ also appears in Connes's prolate heuristic~\cite{Connes26} and in the model and conjecture of~\cite[Prop.~4.1, Conj.~4.2(a)]{rh80} (stated there for real odd characters).
\item \emph{Finite matrices.} Andrade~\cite{Andrade2} and Andrews Jr.~\cite{AndrewsJr} certify simple-even for finite Galerkin matrices $QW^N$; these are not statements about the continuum form.
\item \emph{Classical ancestor.} Yoshida~\cite{Yoshida} combined finitely many pivots with a rigorous high-frequency complement lemma.
\end{itemize}

\subsection{Contributions}
\begin{enumerate}
\item A computer-assisted proof that the ground state of $Q_L$ is simple and even for every $L\in(0,\Lmax]$ (Theorem~\ref{thm:main}). The part we believe to be new is $(0.8,\Lmax]$ (literature search including rStation, 31 October 2026): there, that the ground state is simple and even, and the certified lower bounds for $\lambda_2^\ev$ that prove it. Positivity on this range is not new~\cite{rh45,rh60}. The piece $(0,1.19]$ was released earlier as an unrefereed note~\cite{Botha}; this paper reproduces its mathematical content.
\item A rank-one deflation added to Elston's certifier (Lemma~\ref{lem:deflation}, Corollary~\ref{cor:lambda2}), which turns a certified lowest-eigenvalue bound into a certified \emph{second} even eigenvalue bound at no extra assembly cost. The rStation screen of 31 October 2026 found no bound on $\lambda_2^\ev$; \cite{rh60} proves positivity of both sectors up to $a=2$ without explicit eigenvalue lower bounds.
\item Rigorous upper bounds for the lowest even and odd eigenvalues up to $L=2.25$ from a fast Legendre convolution recurrence with a rigorous Gauss--Legendre bound (\S\ref{sec:upper}). We claim only explicit rigorous values at $L=2.25$, the sharpest we know of there; analytic upper bounds with double-exponential decay are due to Mori~\cite[Thm.~K]{MoriV} and bramble\_trace~\cite[Thm.~2.1]{rh25} (unconditional) and to span\_of\_dilates~\cite[Thm.~1]{rh81} (under RH). No value is claimed to be close to the eigenvalue.
\item Numerical measurements of the weak singular points of the computed minimisers at $|y|=|L-\log n|$ reported by Eberwein~\cite[\S13]{Eberwein} (\S\ref{sec:kinks}); their location is not new.
\end{enumerate}
\emph{Not claimed as new:} the range $(0,0.8]$~\cite{KB} and in particular $(0,0.1406]$~\cite{PL}; the window $L=0.8$ by itself [Zhu, Thm.~6.2]; positivity on this range, and on $(0,2]$~\cite{rh45,rh60}; the $4\pi e^{2L}$ decay scale of the lowest eigenvalue~\cite{Connes26,MoriV,rh80,rh81}; certified upper bounds on the lowest eigenvalue for $L\in[0.5,2]$ [Zhu, Table~3]; the frequency-cutoff reduction and block estimate [Zhu, Thm.~1.1]; the window-compressed prime-norm reduction and its implementation~\cite{Elston}; the Legendre convolution recurrence~\cite{HaleTownsend}, already used for this problem by Liu~\cite{Liu}.

% =====================================================================
\section{Setup and normalisation}\label{sec:setup}

\subsection{Frequency side}
For real $f$ supported in $[-L,L]$ one has $\hat g(t)=F(t)F(-t)=|F(t)|^2$ for real~$t$, and by Plancherel $g(0)=\|f\|^2=\frac1\pi\int_0^\infty|F|^2$. Hence
\begin{gather}\label{eq:freq}
Q_L(f)=[\pole]+\frac1\pi\int_0^\infty\Psi_L(t)|F(t)|^2\,dt,\\
\Psi_L(t)=\re\psi\bigl(\tfrac14+\tfrac{it}2\bigr)-\log\pi-\sum_{\log n<2L}\frac{2\Lambda(n)}{\sqrt n}\cos(t\log n),\notag
\end{gather}
with $[\pole]=2F(\tfrac i2)F(-\tfrac i2)$. Here $F(\tfrac i2)=\int fe^{-x/2}$ is real; for even $f$, $[\pole]=+2F(\tfrac i2)^2$, and for odd $f$, $[\pole]=-2F(\tfrac i2)^2$. This normalisation is that of Zhu [Zhu, \S6]; it was checked numerically against the zero side of the explicit formula (first $400$ zeros).

\subsection{Position side}
Let $\gamma$ denote Euler's constant. For $g=f*\tilde f$,
\begin{multline}\label{eq:P}
Q_L(f)=4\int_0^{2L}g(x)\cosh\tfrac x2\,dx-g(0)\bigl(\log\pi+\gamma+\log(1-e^{-4L})\bigr)\\
+\int_0^{2L}\frac{2\,[e^{-2x}g(0)-e^{-x/2}g(x)]}{1-e^{-2x}}\,dx-2\sum_{\log n<2L}\frac{\Lambda(n)}{\sqrt n}\,g(\log n).
\end{multline}
\emph{Derivation.} For smooth $f$, insert $\psi(s)=-\gamma+\int_0^\infty\frac{e^{-t}-e^{-st}}{1-e^{-t}}\,dt$ into~\eqref{eq:weil} and apply Fourier inversion (Fubini applies since $g$ is smooth). For $x>2L$ only the $g(0)$ part of the integrand survives and integrates in closed form, giving the term $\log(1-e^{-4L})$; the pole term gives $4\int_0^{2L}g\cosh(x/2)$. For a polynomial $f$ on $[-L,L]$ (needed for the trial functions), $g$ is a polynomial on $[0,2L]$ and Lipschitz on~$\R$; apply~\eqref{eq:P} to the mollification $f*\phi_\epsilon$ on the window $L+\epsilon$ (no $\log n$ equals $2L$, and the $\log(1-e^{-4L})$ term absorbs the change of upper limit) and let $\epsilon\to0$: the frequency side converges by dominated convergence since $\Psi_L|F|^2=O(t^{-2}\log t)$, and on the position side $g_\epsilon\to g$ uniformly with a uniform Lipschitz constant, which dominates the integrand $\sim(g(0)-g(x))/(2x)$ near $x=0$.

\subsection{The operator, parity and discreteness}
Let $A_L$ be the self-adjoint operator associated with the closure of $Q_L$ on $C_c^\infty(-L,L)$ in $L^2(-L,L)$. By Lemma~\ref{lem:E} below, $\Psi_L(t)\ge\log(t/2\pi)-1/t-\sum_{\log n<2L}2\Lambda(n)n^{-1/2}\to\infty$, and the pole term is bounded by $C\|f\|^2$; hence the form norm controls $\int\log(2+t)|F|^2$. For functions supported in $[-L,L]$ the set where this quantity and $\|f\|$ are bounded is compact in $L^2$ (high frequencies are uniformly small, and on compact $t$-sets the $F$ are equicontinuous by Paley--Wiener). So $A_L$ has compact resolvent, in agreement with~\cite[Thm.~3.6]{CCM25}. The kernel is real, so we work with real $f$ throughout.

\begin{lemma}[Parity splitting]\label{lem:parity}
Let $(Rf)(x)=f(-x)$. Then $Q_L(Rf)=Q_L(f)$, and $Q_L(f_e+f_o)=Q_L(f_e)+Q_L(f_o)$ for $f_e$ even and $f_o$ odd. Hence $A_L=A_L^{\ev}\oplus A_L^{\od}$.
\end{lemma}
\begin{proof}
$Rf$ has transform $F(-t)$, which leaves $|F(t)|$ ($t$ real) and $F(\tfrac i2)F(-\tfrac i2)$ unchanged. For the bilinear form $B$ of $Q_L$, $B(f_e,f_o)=B(Rf_e,Rf_o)=-B(f_e,f_o)$.
\end{proof}

\begin{definition}
For each parity sector $\sigma\in\{\ev,\od\}$ let $\lambda_1^\sigma(L)\le\lambda_2^\sigma(L)\le\cdots$ be the eigenvalues of $A_L^\sigma$, counted with multiplicity; by compactness of the resolvent they coincide with the min-max values
\[
\lambda_k^\sigma(L)=\inf_{\substack{V\subset\mathcal D^\sigma\\ \dim V=k}}\ \sup_{0\ne f\in V}\frac{Q_L(f)}{\|f\|^2},
\]
where $\mathcal D^\sigma$ is the form domain of the sector. (Indices start at~$1$, as in~\cite{Zhu,Elston}; the earlier note~\cite{Botha} writes $\lambda_0,\lambda_1^{\ev},\lambda_0^{\od}$ for our $\lambda_1^\ev,\lambda_2^\ev,\lambda_1^\od$.)
\end{definition}

\begin{corollary}\label{cor:criterion}
The lowest eigenvalue of $A_L$ is simple with even eigenfunction if and only if
\[
\lambda_1^\ev(L)<\min\bigl(\lambda_2^\ev(L),\lambda_1^\od(L)\bigr).
\]
\end{corollary}

\begin{remark}\label{rem:affine}
The property in Corollary~\ref{cor:criterion} is invariant under $Q\mapsto aQ+b\|\cdot\|^2$ with $a>0$. Conventions that differ from ours by a positive factor (e.g.\ $2\pi$) or a constant shift (e.g.\ the placement of $\log\pi$) therefore do not affect the theorem.
\end{remark}

% =====================================================================
\section{Monotonicity and the grid argument}\label{sec:mono}

\begin{lemma}[Monotonicity in $L$]\label{lem:mono}
For $0<L<L'$, every $k\ge1$ and $\sigma\in\{\ev,\od\}$, $\lambda_k^\sigma(L')\le\lambda_k^\sigma(L)$.
\end{lemma}
\begin{proof}
Extension by zero $\iota:L^2(-L,L)\to L^2(-L',L')$ is an isometry that preserves parity and maps the form domain into the form domain, and $Q_{L'}(\iota f)=Q_L(f)$: the pole and archimedean terms depend only on $F$, and a prime power with $2L\le\log n<2L'$ contributes $g(\log n)=0$ because $g$ is continuous and supported in $[-2L,2L]$ with $g(\pm2L)=0$. The min-max infimum over the larger family of subspaces of $L^2(-L',L')$ (which contains all $\iota V$) can only be smaller.
\end{proof}

\begin{corollary}[Step criterion]\label{cor:step}
Let $L_i<L_{i+1}$, let $U(L_i)\ge\lambda_1^\ev(L_i)$, and let $\ell_2(L_{i+1})\le\lambda_2^\ev(L_{i+1})$ and $\ell^\od(L_{i+1})\le\lambda_1^\od(L_{i+1})$. If
\begin{equation}\label{eq:step}
\min\bigl(\ell_2(L_{i+1}),\ell^\od(L_{i+1})\bigr)>U(L_i),
\end{equation}
then the ground state of $A_L$ is simple and even for every $L\in[L_i,L_{i+1}]$.
\end{corollary}
\begin{proof}
For such $L$, by Lemma~\ref{lem:mono},
\[
\lambda_1^\ev(L)\le\lambda_1^\ev(L_i)\le U(L_i)<\min(\ell_2,\ell^\od)(L_{i+1})\le\min\bigl(\lambda_2^\ev(L),\lambda_1^\od(L)\bigr);
\]
conclude with Corollary~\ref{cor:criterion}.
\end{proof}

We call the quotient $\min(\ell_2,\ell^\od)(L_{i+1})/U(L_i)$ the \emph{step ratio}. Prime-power entries inside a step are harmless, since Lemma~\ref{lem:mono} holds across them.

% =====================================================================
\section{Main theorem}\label{sec:main}

\begin{theorem}\label{thm:main}
For every $L\in(0,\Lmax]$ the lowest eigenvalue of the operator $A_L$ associated with Weil's quadratic form~\eqref{eq:weil} on $L^2(-L,L)$ is simple, and its eigenfunction is even. Quantitatively:
\begin{enumerate}
\item on $(0,0.05]$: $\lambda_2^\ev-\lambda_1^\ev\ge0.4172$ and $\lambda_1^\od-\lambda_1^\ev\ge0.6624$;
\item on each of the $30$ grid steps $[L_i,L_{i+1}]\subset[0.05,1.19]$ of Table~\ref{tab:piece2}: step ratio $\ge1.288$;
\item on each of the $8$ grid steps $[L_i,L_{i+1}]\subset[1.19,\Lmax]$ of Table~\ref{tab:piece3}: step ratio $\ge1.131$.
\end{enumerate}
\end{theorem}

\begin{remark}[Status of release 1.1.0]\label{rem:status}
Items (1) and (2) are the certificate of record of~\cite{Botha}, re-checked by the gluing script (\S\ref{sec:repro}). Item~(3) rests on the $16$ lower-bound certificates and $8$ upper bounds of Table~\ref{tab:piece3}, all complete. The goal of this work is $L\in(0,1.5]$: the extension to $1.5$ is in progress, with the even sector at $1.4255$ and both sectors at $1.468$, $1.4885$ and $1.5$ still running. The later step $[1.4255,1.447]$ already passes (Table~\ref{tab:piece3}), but it is not part of the theorem, because the step $[1.4035,1.4255]$ before it is not yet complete. At the single window $L=1.19$ the Piece~III machinery also confirms simple-even with a second, independent code path (\S\ref{sec:results3}). Piece~III beyond $L=1.19$ has not been audited (\S\ref{sec:limitations}).
\end{remark}

The numbers in (1)--(3) are outward-rounded lower ends. The proof is computer-assisted and uses Arb ball arithmetic~\cite{Johansson,FLINT}; it is given in \S\ref{sec:proof}.

% =====================================================================
\section{Proof}\label{sec:proof}

\subsection{Piece I: small windows $L\in(0,0.05]$}\label{sec:piece1}
Here $L<\tfrac12\log2$, so no prime powers occur. Put $f(x)=L^{-1/2}f_1(x/L)$ with $\|f_1\|_{L^2(-1,1)}=1$ and $F_1=\hat f_1$. The substitution $t=\tau/L$ gives
\[
Q_L(f)=\log\tfrac1L+R_L(f_1),\qquad R_L(f_1)=P_L(f_1)+\frac1\pi\int_0^\infty m_L(\tau)|F_1(\tau)|^2\,d\tau,
\]
with $m_L(\tau)=\re\psi(\tfrac14+\tfrac{i\tau}{2L})-\log\pi+\log L$, $P_L(f_1)=2L\bigl(\int_{-1}^1f_1e^{-Ls/2}ds\bigr)^2$ for even $f_1$ and $P_L(f_1)=-2L\bigl(\int_{-1}^1f_1\sinh(Ls/2)ds\bigr)^2$ for odd $f_1$. The shift $\log(1/L)$ is common to all $f_1$, so the gaps of $Q_L$ equal those of $R_L$, sector by sector, and $R_L$ acts on the fixed space $L^2(-1,1)$. With $u=\tau/(2L)$,
\[
\partial_Lm_L=\tfrac1L h(u),\qquad h(u)=1-\re\bigl(iu\,\psi'(\tfrac14+iu)\bigr).
\]

\begin{lemma}\label{lem:h}
For $u>0$, $|h(u)|\le0.8125/u^2$.
\end{lemma}
\begin{proof}
From $\psi'(w)=\sum_{n\ge0}(n+w)^{-2}$ one gets $h(u)=1-\sum_{n\ge0}\varphi(n+\tfrac14)$ with $\varphi(a)=2au^2/(a^2+u^2)^2$ and $\int_0^\infty\varphi=1$. The sum is the midpoint rule for $\int_{-1/4}^\infty\varphi$ with unit cells. The midpoint Peano kernel is bounded by $1/8$ and $\int_\R|\varphi''|=6/u^2$, so the quadrature error is at most $\tfrac34u^{-2}$; and $|\int_{-1/4}^0\varphi|\le\int_{-1/4}^02|a|u^{-2}da=\tfrac1{16}u^{-2}$.
\end{proof}

We bound $\int_0^\infty|h|$ by an upper Riemann sum on $[0,40]$ with cells of width $1/200$, enclosing $\sup|h|$ on each cell by evaluating $h$ on a ball covering the cell, plus $0.8125/40$ from Lemma~\ref{lem:h}. At $768$ bits this gives $\int|h|\le0.593979$. Since $|F_1|^2\le(\int|f_1|)^2\le2$, the multiplier part of $\frac{d}{dL}R_L(f_1)$ is at most $\frac4\pi\int|h|$. For the even pole term, with $J=\int f_1e^{-Ls/2}$, $\frac{d}{dL}(2LJ^2)=2J^2+4LJJ'$ and Cauchy--Schwarz gives $J^2\le2\sinh L/L$, $|J'|\le\tfrac12(\tfrac23e^L)^{1/2}$; for the odd pole term, with $K=\int f_1\sinh(Ls/2)$, $K^2\le2\sinh^2(L/2)$ and $|K'|\le\tfrac12\cosh(L/2)\sqrt{2/3}$. All bounds increase in $L$; evaluated at $L_0=0.05$ they give, uniformly for $L\le L_0$,
\[
C_\ev=\tfrac4\pi\textstyle\int|h|+2J^2+4L_0|J||J'|\le4.876362,\qquad C_\od\le0.761666.
\]
The forms $R_L$ share the core $C_c^\infty(-1,1)$ and $|R_L(f_1)-R_{L'}(f_1)|\le C_\sigma|L-L'|\|f_1\|^2$ in sector $\sigma$; by min-max every sector eigenvalue of $R_L$ is $C_\sigma$-Lipschitz in $L$. With the certified values $\ell_1,U,\ell_2,\ell^\od$ at $L_0=0.05$ (Table~\ref{tab:piece2}), for $L\le L_0$,
\begin{gather*}
\lambda_2^\ev(L)-\lambda_1^\ev(L)\ge\ell_2-U-2C_\ev L_0\ge0.4172,\\
\lambda_1^\od(L)-\lambda_1^\ev(L)\ge\ell^\od-U-(C_\ev+C_\od)L_0\ge0.6624.
\end{gather*}
Both are positive, which proves item~(1) of Theorem~\ref{thm:main}. \qed

\begin{remark}
At $53$ bits the same procedure gives the smaller enclosures $\int|h|\le0.592773$, $C_\ev\le4.87484$, $C_\od\le0.760130$; at $128$ bits $\int|h|\le0.593583$. All are upper bounds under Arb's containment guarantee and we use the largest. Why the higher-precision upper end is larger is not explained (\S\ref{sec:limitations}).
\end{remark}

\subsection{Per-window certificates on $[0.05,1.19]$}\label{sec:piece2cert}
Fix $L$ and a parity. We use the orthonormal Legendre basis $e_n(x)=c_nP_n(x/L)$, $c_n=\sqrt{(2n+1)/(2L)}$, with $n=2m$ (even) or $n=2m+1$ (odd). Up to a unit factor common to the sector, $\hat e_n(t)=c_n\,2L\,(-1)^mj_n(Lt)$ with $j_n$ the spherical Bessel function, and the pole vector is $p_m=\sqrt2\,c_n\,2L\,i_n(L/2)$, so that $[\pole]=\pm(p\cdot a)^2$ for $f=\sum a_me_{n(m)}$.

\begin{lemma}[Binet bound]\label{lem:B}
For $\re w>0$, $\bigl|\psi(w)-\log w+\tfrac1{2w}\bigr|\le\dfrac1{12|w|\re w}$.
\end{lemma}
\begin{proof}
Binet's formula reads $\psi(w)=\log w-\frac1{2w}-\int_0^\infty b(x)e^{-wx}dx$ with $b(x)=\frac1{e^x-1}-\frac1x+\frac12$, $b(0)=0$. With $y=x/2$, $b'(x)=\tfrac14(y^{-2}-\sinh^{-2}y)\ge0$. For $y<\sqrt3$ the inequality $b'\le\tfrac1{12}$ reduces to $\sinh^2y\le3y^2/(3-y^2)$, which holds coefficientwise ($2^{2k-1}/(2k)!\le3^{1-k}$, with equality at $k=1,2$ and term ratio $4/((2k+1)(2k+2))\le\tfrac13$); for $y\ge\sqrt3$ it is trivial. Integrating by parts, $\int be^{-wx}=w^{-1}\int b'e^{-wx}$, of modulus at most $1/(12|w|\re w)$.
\end{proof}

\begin{lemma}[Envelope; {[Zhu, Lemma 3.1]}]\label{lem:E}
For $w=\tfrac14+\tfrac{it}2$ and $t\ge3/2$, $\re\psi(w)\ge\log\tfrac t2-\tfrac2{3t}-\tfrac1{2t^2}\ge\log\tfrac t2-\tfrac1t$.
\end{lemma}
\begin{proof}
Lemma~\ref{lem:B} with $|w|\ge t/2$ and $\re(1/(2w))\le1/(2t^2)$; the second inequality is equivalent to $t\ge3/2$.
\end{proof}

Let $A_L^{\sup}=\sum_{\log n<2L}2\Lambda(n)/\sqrt n$. The prime comb is $\ge-A_L^{\sup}$ and the envelope increases in $t$, so for $t\ge\Tsh$
\begin{equation}\label{eq:beta}
\Psi_L(t)\ge\beta:=\log\frac{\Tsh}{2\pi}-\frac1{\Tsh}-A_L^{\sup}.
\end{equation}

\begin{lemma}[Frequency reduction; {[Zhu, Thm.~1.1]}]\label{lem:R}
Assume $\Psi_L\ge\beta>0$ on $[\Tsh,\infty)$ and let $\mathcal R(f)=[\pole]+\frac1\pi\int_0^{\Tsh}(\Psi_L-\beta)|F|^2dt+\beta\|f\|^2$. Then $Q_L\ge\mathcal R$ as forms, hence $\lambda_k^\sigma(Q_L)\ge\lambda_k^\sigma(\mathcal R)$ for all $k$ and $\sigma$.
\end{lemma}
\begin{proof}
Write $\Psi_L=(\Psi_L-\beta)+\beta$, use $\frac1\pi\int_0^\infty|F|^2=\|f\|^2$, drop the nonnegative integral over $[\Tsh,\infty)$ and apply min-max.
\end{proof}

\begin{lemma}[Monotonicity of $\re\psi$]\label{lem:M}
For $s>0$, $\re\psi(s+iy)=\psi(s)+\sum_{n\ge0}\frac{y^2}{(n+s)((n+s)^2+y^2)}$ is increasing in $|y|$. Consequently $|\Psi_L-\beta|\le M_\psi:=\max(|\re\psi(\tfrac14)|,|\re\psi(\tfrac14+\tfrac{i\Tsh}2)|)+\log\pi+A_L^{\sup}+\beta$ on $[0,\Tsh]$.
\end{lemma}
\begin{proof}
Take real parts in $\psi(w)=-\gamma+\sum_n\bigl(\frac1{n+1}-\frac1{n+w}\bigr)$.
\end{proof}

\begin{lemma}[Bessel bounds]\label{lem:J}
$|j_n(x)|\le\dfrac{x^n}{(2n+1)!!}$ and $0\le i_n(x)\le\dfrac{x^n}{(2n+1)!!}e^{x^2/2}$ for $x\ge0$.
\end{lemma}
\begin{proof}
Use Poisson's integral and the beta integral,
\[
j_n(x)=\frac{x^n}{2^{n+1}n!}\int_{-1}^1\cos(xs)(1-s^2)^n\,ds,\qquad\int_{-1}^1(1-s^2)^n\,ds=\frac{2^{2n+1}(n!)^2}{(2n+1)!};
\]
for $i_n$ replace $\cos$ by $\cosh\le e^{x^2/2}$.
\end{proof}

\begin{lemma}[Gauss--Legendre and Bernstein; {\cite[Thm.~19.3]{Trefethen}}]\label{lem:G}
If $g$ is analytic in the Bernstein ellipse $E_\rho$ with $|g|\le M$ there, then the $n$-point Gauss--Legendre rule on $[-1,1]$ has error at most $\frac{64}{15}M\rho^{-2n+2}/(\rho^2-1)$. A polynomial of degree $d$ bounded by $M$ on $[-1,1]$ is bounded by $M\rho^d$ on $E_\rho$.
\end{lemma}
\noindent The cited theorem is stated with $\rho^{-2n}$ for the $(n+1)$-point rule; the form above is the same statement for $n$ points.

\begin{remark}[Correction of a Gauss--Legendre constant]\label{rem:erratum}
The original runs of the window certificates applied $\rho^{-2n}$ to the $n$-point rule, so quadrature balls were too small by $\rho^2$ (about $4.33$ for the head panels, $9$ for the Rayleigh outer rule). The error was found in internal review and fixed in the code. The original certificate was first corrected by a Weyl shift of every lower bound by $\delta=N(\rho^2-1)q_{\mathrm{old}}$ ($\delta\le6.4\e{-35}$ for $L\le0.8$, $\le1.6\e{-33}$ for $0.84\le L\le0.97$, $\le3.2\e{-32}$ for $0.98\le L\le1$) and fresh upper bounds; that correction assumed unchanged window parameters between code revisions. It has since been superseded by a clean end-to-end rerun with the corrected code, which is the certificate of record used here and does not depend on the correction.
\end{remark}

\emph{Head matrix.} We compute $C_{jk}=\frac1\pi\int_0^{\Tsh}(\Psi_L-\beta)\hat e_j\hat e_k\,dt$ by $64$-point Gauss--Legendre on each unit panel in Arb. The ellipse has half-height $B=2/5$, $\rho=2B+\sqrt{4B^2+1}$. For complex $t$ the continuation of $\re\psi(\tfrac14+\tfrac{it}2)$ is $\tfrac12(\psi(w_+)+\psi(w_-))$, $w_\pm=\tfrac14\pm\tfrac{it}2$; on the strip $\re w_\pm\ge\tfrac14-\tfrac B2=0.05$, and Lemma~\ref{lem:B} bounds $\psi(w_\pm)$. On the strip $|\hat e_j|\le\sqrt{2L}e^{LB}$ and the comb is bounded by $\sum2\Lambda(n)n^{-1/2}\cosh(B\log n)$. The per-entry error is $\Tsh\cdot\frac{32}{15}M\rho^{-126}/(\rho^2-1)$ and is added as a ball. The head is $A=C\pm pp^T+\beta I+(\text{quadrature ball})$, with $+$ in the even and $-$ in the odd sector.

\emph{Tails.} For tail orders $n\ge n_0$ put $X=L\Tsh$; then $|\hat e_n|\le\sigma_n:=c_n2LX^n/(2n+1)!!$ on $[0,\Tsh]$ and $\sigma_{n+2}/\sigma_n\le r=\sqrt{\frac{2n_0+5}{2n_0+1}}\frac{X^2}{(2n_0+3)(2n_0+5)}<1$ (asserted), so $\sum\sigma^2\le\sigma_{n_0}^2/(1-r^2)$. The tail of the pole vector satisfies $\|p_{\mathrm{tail}}\|\le p_{\mathrm{tail}}:=\sqrt2c_{n_0}2Ly^{n_0}/(2n_0+1)!!\cdot e^{y^2/2}/\sqrt{1-y^2}$, $y=L/2$. Put
\[
e_D=\frac{M_\psi}\pi\Tsh\sum\sigma^2+[\od]\,p_{\mathrm{tail}}^2,\qquad
e_B=\frac{M_\psi}\pi\Tsh\sqrt{2L}\sqrt N\Bigl(\sum\sigma^2\Bigr)^{1/2}+2\|p_{\mathrm{head}}\|p_{\mathrm{tail}}.
\]
Then the tail block satisfies $D\ge(\beta-e_D)I$ (in the even sector the tail pole term $+pp^T\ge0$ is dropped) and the coupling block satisfies $\|X\|\le e_B$ by Hilbert--Schmidt bounds.

\begin{lemma}[Block bound; {[Zhu, Thm.~1.1, \S6.2]}]\label{lem:S}
Let $\mathcal R=\begin{pmatrix}A&X\\X^T&D\end{pmatrix}$ with $D\ge(\beta-e_D)I$ and $\|X\|\le e_B$. Then $\lambda_k(\mathcal R)\ge\min(\lambda_k(A),\beta-e_D)-e_B$ for every $k$.
\end{lemma}
\begin{proof}
$\bigl\|\begin{pmatrix}0&X\\X^T&0\end{pmatrix}\bigr\|=\|X\|$, so $\mathcal R\ge\operatorname{diag}(A,D)-e_B$. The $k$-th min-max value of $\operatorname{diag}(A,D)$ is at least $\min(\lambda_k(A),\lambda_1(D))$; conclude with Weyl's inequality.
\end{proof}

\emph{Verified Cholesky.} To show $A-\mu I\succeq0$ for a ball matrix $A$, set $d=\mu/100$, compute an approximate Cholesky factor $R$ of $\operatorname{mid}(A)-(\mu+d)I$ in midpoint arithmetic, compute $E=A-(\mu+d)I-RR^T$ rigorously in ball arithmetic, and check $\|E\|_F<d$; then $A-\mu I=RR^T+E+dI\succeq0$ for every matrix in the ball (cf.~\cite{Rump}). Eigenvalue estimates from inverse iteration on midpoints are used only to choose $\mu$ and the deflation vector and carry no rigour. With Lemmas~\ref{lem:R} and~\ref{lem:S} this gives $\lambda_1^\ev\ge\min(\mu_1,\beta-e_D)-e_B$, $\lambda_1^\od\ge\min(\mu^\od,\beta-e_D^\od)-e_B^\od$, and $\lambda_2^\ev\ge\min(\mu_2,\beta-e_D)-e_B$ where $\mu_2$ is certified for $A+vv^T$ (Lemma~\ref{lem:deflation} below).

\emph{Parameters.} The cutoff is $\Tsh=\max(200,5\lceil80/(5L)\rceil,5\lceil2\pi e^{A_L^{\sup}+1/2}/5\rceil)$ and the head has $N=\lceil eL\Tsh/4\rceil+42$ modes per sector. For $L>1$ (``sharp cutoff'') we use instead $\Tsh=2450$ ($L=1.03$), $4000$ ($L=1.06,1.09$), $5550$ ($L=1.12$--$1.19$), the exact dyadic floor $\beta=3/16$, $112$ Gauss nodes per unit panel and $N=\lceil eL\Tsh/4\rceil+80$. Lemma~\ref{lem:E} is applied above the crude cutoff $T_2$ (asserting $\log(T_2/2\pi)-1/T_2-A_L^{\sup}\ge3/16$), and $\Psi_L>\beta$ on $[\Tsh,T_2]$ is certified by Arb enclosures of $\Psi_L$ on cells of width $0.01$ with bisection (minimum width $10^{-9}$). Together $\Psi_L\ge\beta$ on $[\Tsh,\infty)$, which is all Lemma~\ref{lem:R} needs. This device is a refinement inside Zhu's reduction; we make no novelty claim for it. Validation at $L=1$: the sharp path gives $4.978\e{-30}\le\lambda_1^\ev\le9.406\e{-30}$, $\lambda_2^\ev\ge1.808\e{-23}$, $\lambda_1^\od\ge1.255\e{-26}$, overlapping the crude-cutoff values of Table~\ref{tab:piece2}.

\emph{Spherical Bessel values for $L\ge0.98$.} Orders $n\le n_1=\lfloor x\rfloor$ use upward recurrence in ball arithmetic from Arb seeds every $400$ orders; orders $n>n_1$ use the ratios $r_k=j_k/j_{k-1}$ from the backward continued fraction $r_k=x/(2k+1-xr_{k+1})$, started from $[0,1]$ and anchored at $j_{n_1}$.

\begin{lemma}[Ratio enclosure]
For $x>0$ and $k\ge\lfloor x\rfloor+1$, every true ratio $r_k$ lies in $(0,1]$; hence the interval iteration started from $[0,1]$ encloses it.
\end{lemma}
\begin{proof}
If $2k+1\ge2x$, the map $\tau\mapsto x/(2k+1-x\tau)$ sends $[0,1]$ into $[x/(2k+1),x/(2k+1-x)]\subset[0,1]$. By downward induction from large $K$ (where $j_K(x)>0$ and $r_K\in[0,1]$): if $2k+1\ge2x$ and $r_{k+1}\in[0,1]$, then $j_{k-1}/j_k=(2k+1)/x-r_{k+1}\ge1$, so $r_k\in(0,1]$ and $j_{k-1}>0$. The condition holds for all $k\ge\lfloor x\rfloor+1$.
\end{proof}
The enclosures need not be sharp at the top orders unless the start index is well above $n_{\max}$~\cite{Gautschi}; sharpness only affects the size of the radii.

\emph{Upper bound $U$.} The trial function is $f=\sum_{m<m^*}v_me_{2m}$, a polynomial on $[-L,L]$, with $v$ an approximate ground vector (exact dyadic entries). In $s=x/L$, $g(Ls)=L\int_{-1}^{1-s}f(y)f(y+s)dy$ is computed exactly by $(D+1)$-point Gauss--Legendre ($D=\deg f$); the integral in~\eqref{eq:P} by $450$-point Gauss--Legendre on $[0,2]$ with error bounded by Lemma~\ref{lem:G} ($\rho=3$, exponent $\rho^{-898}$, ellipse maximum enclosed by $2048$ ball arcs with $|1-e^{-2Lz}|>0$ asserted on each arc; the zeros $i\pi k/L$, $k\neq0$, lie outside the ellipse); the other terms of~\eqref{eq:P} exactly. Then $U\ge Q_L(f)/\|f\|^2\ge\lambda_1^\ev(L)$. A polynomial on $[-L,L]$ has $|\hat f|^2=O(t^{-2})$ and lies in the form domain $\{f:\int\log(2+t)|\hat f|^2<\infty\}$, and in the closure of $C_c^\infty(-L,L)$ there, since the weight is weaker than any $H^s$, $s<\tfrac12$.

\subsection{Piece II: the grid on $[0.05,1.19]$}\label{sec:piece2}
The grid is
\begin{gather*}
0.05,\ 0.1,\ 0.2,\ 0.3,\ 0.35,\ 0.4,\ 0.45,\ 0.5,\ 0.55,\ 0.6,\ 0.64,\ 0.68,\ 0.72,\ 0.76,\ 0.8,\ 0.84,\\
0.87,\ 0.9,\ 0.93,\ 0.96,\ 0.97,\ 0.98,\ 0.99,\ 1.0,\ 1.03,\ 1.06,\ 1.09,\ 1.12,\ 1.15,\ 1.18,\ 1.19.
\end{gather*}
For each of the $30$ steps the criterion~\eqref{eq:step} holds (Table~\ref{tab:piece2}); the smallest ratios are $1.288$ on $[0.05,0.1]$, $1.320$ on $[0.1,0.2]$ and, past $L=1$, $2.148$ on $[1.15,1.18]$. By Corollary~\ref{cor:step} this proves item~(2) of Theorem~\ref{thm:main}. \qed

\begin{remark}[Comparisons and dyadic windows]
(a) The gluing script compares balls with Python's built-in \texttt{min}, which is not a rigorous minimum when balls overlap; here the two lower bounds are disjoint at every window (relative separation $\ge1.021$, radii $<5\e{-20}$), and the step ratios in Table~\ref{tab:piece2} were recomputed with explicit endpoint comparisons. (b) Each window is certified at the dyadic number $\tilde L_i\ge L_i$ with $\tilde L_i<L_i+2^{-p}$, $p\ge768$, for all four bounds of the window; the covering is then $(0,\tilde L_1]\cup\bigcup_i[\tilde L_i,\tilde L_{i+1}]\supset(0,1.19]$.
\end{remark}

\subsection{Piece III: the grid on $(1.19,\Lmax]$}\label{sec:piece3}
Beyond $L\approx1.2$ the route of \S\ref{sec:piece2cert} is impractical: the cutoff $\Tsh\gtrsim2\pi e^{A_L^{\sup}}$ forced by the pointwise bound~\eqref{eq:beta} is about $2\cdot10^5$ at $L=1.3$ and $5\cdot10^6$ at $L=1.5$, and Zhu [Lemma 3.2] shows that $A_L^{\sup}$ is optimal for any pointwise bound. Elston~\cite{Elston} replaces the pointwise bound by an operator bound for the prime comb on a slightly widened window. We state his reduction as a cited proposition; we use his code unchanged except for the deflation option described below.

\emph{Notation.} Let $a=L$ be the window, $\delta>0$, $b=a+\delta$, $\beta_K>0$, and $\Tsh>\beta_K/\delta$ with $\Omega:=\Tsh-\beta_K/\delta$ and $\Omega\delta\ge1$. Let $W(u)=I_0(\beta_K\sqrt{1-u^2/\delta^2})/I_0(\beta_K)$ on $[-\delta,\delta]$ (Kaiser window, $0$ outside), $\kappa(u)=\frac{\sin\Omega u}{\pi u}W(u)$ and $m=1-\hat\kappa$. Let $H(t)=\re\psi(\tfrac14+\tfrac{it}2)-\log\pi$, $P(t)=\sum_{\log n<2a}\frac{2\Lambda(n)}{\sqrt n}\cos(t\log n)$, so $\Psi_a=H-P$, and $A=\sum_{\log n<2a}2\Lambda(n)/\sqrt n$. Let $C_b=\sum_{\log n<2a}\Lambda(n)n^{-1/2}(\tau_{\log n}+\tau_{-\log n})$, compressed to $L^2(-b,b)$, where $\tau_s$ is translation by~$s$.

\begin{proposition}[Elston~\cite{Elston}]\label{prop:elston}
Assume
\begin{enumerate}
\item[(H1)] $\bar\mu\ge\sup\{\langle h,C_bh\rangle/\|h\|^2:0\ne h\in L^2(-b,b)\}$;
\item[(H2)] $\eps\ge\sup_{t\ge\Tsh}|\hat\kappa(t)|$ (in~\cite{Elston}, $\eps=2/(\pi\beta_KI_0(\beta_K))$);
\item[(H3)] $0<\beta'\le\log\frac{\Tsh}{2\pi}-\frac1{\Tsh}-\bar\mu-(2\eps+\eps^2)(A+\bar\mu)$.
\end{enumerate}
Put $\Phi=\Psi_a+m^2(P-\bar\mu)$ and
\[
R'(f)=[\pole]+\frac1\pi\int_0^{\Tsh}(\Phi-\beta')|F|^2\,dt+\beta'\|f\|^2.
\]
Then:
\begin{enumerate}
\item[(a)] $Q_a\ge R'$ as forms on each parity sector.
\item[(b)] Let $A_N$ be the matrix of $R'$ on the first $N$ Legendre modes of the sector, and let $\eps_D,\eps_B$ be the tail constants of Zhu's reduction computed with $G\ge\sup_{[0,\Tsh]}|\Phi-\beta'|$ (Bessel tail bounds as in Lemma~\ref{lem:J}; $\eps_D$ by a Gershgorin bound on the tail rows, $\eps_B$ by a Schur test). Then the tail block of $R'$ is $\succeq(\beta'-\eps_D)I$ and the coupling block has norm $\le\eps_B$.
\item[(c)] The certifier encloses every entry of $A_N$ in a ball (Gauss--Legendre nodes and weights, $\Psi$ via the complex digamma function, $m$ by rigorous integration of the Kaiser main lobe and a mean-value bound for the side lobes, $j_n$ by recurrence and ratio enclosure), adds a spectral quadrature debit $c_{\mathrm{err}}$ from a Bernstein-ellipse bound on panels of width $1/4$, and certifies a lower bound $\lambda_1(A_N)\ge\mu'-\rho_{\mathrm{res}}-N r_{\max}-c_{\mathrm{err}}$, where $\mu'$ is a dyadic shift, $\rho_{\mathrm{res}}$ the $\infty$-norm of an exactly computed integer Cholesky residual divided by the scale $2^{2P}$ (plus $N2^{-2P}$ for the floors), and $r_{\max}$ the largest entry radius.
\end{enumerate}
\end{proposition}
\begin{proof}[Proof of (a), following~\cite{Elston}]
Put $h=f-f*\kappa$; then $\operatorname{supp}h\subset[-b,b]$ and $\hat h=mF$. By (H1), $\frac1\pi\int_0^\infty Pm^2|F|^2=\langle h,C_bh\rangle\le\bar\mu\|h\|^2=\frac1\pi\int_0^\infty\bar\mu m^2|F|^2$, so $\frac1\pi\int m^2(\bar\mu-P)|F|^2\ge0$ and $Q_a(f)\ge[\pole]+\frac1\pi\int_0^\infty\Phi|F|^2$. For $t\ge\Tsh$, $\Phi=H-\bar\mu-(1-m^2)(P+\bar\mu)$ with $|1-m^2|\le2\eps+\eps^2$ by (H2), so Lemma~\ref{lem:E} and (H3) give $\Phi\ge\beta'$; dropping $\frac1\pi\int_{\Tsh}^\infty(\Phi-\beta')|F|^2\ge0$ gives $Q_a\ge R'$.
\end{proof}
Elston states (b)--(c) for the lowest eigenvalue; in~\cite{Elston}, $\bar\mu$ is certified by a weighted Schur test on $4000$--$8000$ cells with a Collatz--Wielandt weight, and the finite certificate is an exact integer Gram residual in the spirit of~\cite{Rump}. An independent audit of the code at $L=1.19$ (documented in~\cite{Elston}) found and fixed two implementation issues; the code used here is the post-audit version (upstream commit \file{3cc4cb7b}).

\begin{lemma}[Rank-one deflation]\label{lem:deflation}
Let $A$ be a real symmetric $N\times N$ matrix and $w\in\R^N$. Then $\lambda_2(A)\ge\lambda_1(A+ww^T)$. In particular, if $A+ww^T-\mu I\succeq0$ then $A$ has at most one eigenvalue below $\mu$.
\end{lemma}
\begin{proof}
By Courant--Fischer, $\lambda_2(A)=\max_{y\ne0}\min\{x^TAx:\|x\|=1,\,x\perp y\}\ge\min\{x^TAx:\|x\|=1,\,x\perp w\}=\min\{x^T(A+ww^T)x:\|x\|=1,\,x\perp w\}\ge\lambda_1(A+ww^T)$. (For $w=0$ the claim is trivial.)
\end{proof}

\begin{corollary}[Second even eigenvalue]\label{cor:lambda2}
Under the hypotheses of Proposition~\ref{prop:elston}, in the even sector and for any $w\in\R^N$,
\[
\lambda_2^\ev(L)\ \ge\ \lambda_2(R'^{\,\ev})\ \ge\ \min\bigl(\lambda_1(A_N+ww^T),\ \beta'-\eps_D\bigr)-\eps_B .
\]
\end{corollary}
\begin{proof}
By Lemma~\ref{lem:deflation}, $\lambda_2(A_N)\ge\lambda_1(A_N+ww^T)$. Lemma~\ref{lem:S} with $k=2$ (Weyl's inequality applied to the off-diagonal blocks, with the bounds of Proposition~\ref{prop:elston}(b)) gives $\lambda_2(R')\ge\min(\lambda_2(A_N),\beta'-\eps_D)-\eps_B$. Since $Q_L\ge R'$ as forms (Proposition~\ref{prop:elston}(a)), min-max gives $\lambda_2^\ev(L)\ge\lambda_2(R'^{\,\ev})$.
\end{proof}

\emph{Implementation of the deflation} (option \texttt{--deflate}; our only change to~\cite{Elston}). After the ball matrix $A_N$ is assembled, a vector $v$ approximating the unit lowest eigenvector of $\operatorname{mid}(A_N)$ is computed by $12$ steps of inverse iteration; its entries are rounded to $20$ significant decimal digits and converted to Arb balls. The matrix $A_N+vv^T$ is formed in ball arithmetic, so it encloses $A_N+ww^T$ for the exact midpoint vector $w$ of $v$, and the rounding of $v$ enters the largest entry radius $r_{\max}$ exactly like the Gram radii. The certifier then runs unchanged on $A_N+vv^T$, and Proposition~\ref{prop:elston}(c) yields a certified lower bound for $\lambda_1(A_N+ww^T)$. The choice of $v$ affects only sharpness, not rigour.

\emph{One form, two codes.} The step test~\eqref{eq:step} compares lower bounds from Elston's certifier with upper bounds from our code, so both must refer to the same form with the same norm. Both implement~\eqref{eq:weil} in Zhu's normalisation ($\|f\|^2=\frac1\pi\int_0^\infty|F|^2$ for real $f$, pole term $+2F(\tfrac i2)^2$ in the even and $-2|F(\tfrac i2)|^2$ in the odd sector). As a consistency check, at $L=1.19$ Elston's certified even lower bound $6.8131\e{-48}$, our Piece~II bracket $[7.8\e{-48},1.312\e{-47}]$, the Legendre--Galerkin bound $U(1.19)\le8.712\e{-48}$ ($K=640$, below) and the non-rigorous split-window value $8.676\e{-48}$ (basis of \S\ref{sec:kinks}) are mutually compatible; in the odd sector Elston's certified lower bound $3.9196\e{-44}$ lies below our rigorous Galerkin upper bound $4.8461\e{-44}$ ($K=256$). The proof does not rely on these numerical agreements; it relies only on both codes evaluating~\eqref{eq:weil}.

\emph{Upper bounds $U(L_i)$.} At every grid point $L_i$, including $L_i=1.19$, $U(L_i)$ is the outward-rounded upper endpoint of a Rayleigh quotient $v^TQ_Kv/v^Tv$ of the even Legendre--Galerkin matrix $Q_K$ ($K=640$ even modes, $3000$ bits) with an exact dyadic vector~$v$, computed in ball arithmetic with the rigorous quadrature bound of \S\ref{sec:upper} (Proposition~\ref{prop:upper}). At $L=1.19$ this gives $U(1.19)\le8.712\e{-48}$, sharper than the $1.312\e{-47}$ of Piece~II (Table~\ref{tab:piece2}). The sharper value is needed: with $1.312\e{-47}$ the first step ratio would be $0.93$. Both are valid upper bounds, and the gluing script uses the smaller one at each window.

\emph{Grid.} The planned grid to $L=1.5$ is
\begin{gather*}
1.19,\ 1.2205,\ 1.25,\ 1.2785,\ 1.3055,\ 1.3315,\ 1.3565,\\
1.3805,\ 1.4035,\ 1.4255,\ 1.447,\ 1.468,\ 1.4885,\ 1.5 .
\end{gather*}
It was chosen from non-rigorous Galerkin estimates so that $0.95\min(\lambda_2^\ev,\lambda_1^\od)(L_{i+1})\ge1.25\cdot1.03\cdot\lambda_1^\ev(L_i)$. In this release the windows $1.2205,\dots,\Lmax$ are complete in both sectors, as are the odd sector at $1.4255$ and both sectors at $1.447$; the other runs are in progress. The parameters actually used are in Table~\ref{tab:grid3} and in the field \texttt{args} of every certificate file. They follow a design rule (\file{code/gap/params.py}) that plays no role in the proof: $\Omega=\Tsh-\beta_K/\delta$ is interpolated log-linearly from Elston's runs at $L=1.19,1.3,1.495$ and inflated by $3\%$; $N\approx0.8L(\Omega+64/\delta)+20$; the number of Gauss points per panel is raised so that the quadrature debit is at most $10^{-3}$ times the target eigenvalue (computed in the certifier's plan mode); and $\beta_K$ is the smallest value for which $N$ times the side-lobe entry radius stays below $10^{-6}$ times the target eigenvalue. The last rule was added after a first pass with $\beta_K=64$ at every window gave a negative odd bound at $1.2785$; the runs from that pass that gave useful bounds ($1.2205$, $1.25$, and the even sector at $1.2785$) are kept. Grid values are passed to the certifier as decimal strings, converted to Arb balls that contain the exact decimal; the Schur window $b$ is rounded up, which is admissible because the supremum in (H1) is monotone in~$b$.

\emph{Conclusion of Piece III.} Each of the $8$ steps from $1.19$ to $\Lmax$ satisfies~\eqref{eq:step} (Table~\ref{tab:piece3}); the smallest ratio is $1.131$, on $[1.2205,1.25]$, and the odd sector binds at every window. By Corollary~\ref{cor:step} this proves item~(3) of Theorem~\ref{thm:main}. \qed

The step $[1.4255,1.447]$ also satisfies~\eqref{eq:step} (ratio $\ge1.379$), so the ground state is simple and even on that interval too. We report it as a data point and not as part of Theorem~\ref{thm:main}, since $[1.4035,1.4255]$ is not yet covered. Nothing is claimed for $L\in(\Lmax,1.4255)$ or $L>1.447$.

% =====================================================================
\section{Results}\label{sec:results}

\subsection{Pieces I and II}
Table~\ref{tab:piece2} lists the certificate of record on $[0.05,1.19]$. Lower bounds are rounded down and upper bounds up, to four significant digits; the step ratio in row $L_{i+1}$ is $\min(\ell_2,\ell^\od)(L_{i+1})/U(L_i)$, lower end floored. The odd sector binds at every window from $L=0.1$ on. Piece~I on this certificate: $\lambda_2^\ev-\lambda_1^\ev\ge0.4172$ and $\lambda_1^\od-\lambda_1^\ev\ge0.6624$ on $(0,0.05]$.

\begin{table}[ht]
\centering\footnotesize
\caption{Certificate of record on $[0.05,1.19]$. $\Tsh$ and $N$ are the cutoff and the head size per sector. At $L=1$ the crude cutoff is used (the sharp-path validation is in \S\ref{sec:piece2cert}).}\label{tab:piece2}
\begin{tabular}{rrrllllr}
\toprule
$L$ & $\Tsh$ & $N$ & $\lambda_1^\ev\ge$ & $\lambda_1^\ev\le U$ & $\lambda_2^\ev\ge$ & $\lambda_1^\od\ge$ & step ratio $\ge$\\
\midrule
0.05 & 1600 & 97 & 0.9351 & 0.9827 & 1.887 & 1.926 & (Piece I)\\
0.1 & 800 & 97 & 0.4348 & 0.4572 & 1.387 & 1.266 & 1.288\\
0.2 & 400 & 97 & 0.09014 & 0.09499 & 1.046 & 0.6037 & 1.320\\
0.3 & 270 & 98 & 7.143e-3 & 7.587e-3 & 0.7437 & 0.2110 & 2.222\\
0.35 & 230 & 97 & 1.110e-3 & 1.204e-3 & 0.6153 & 0.06086 & 8.022\\
0.4 & 200 & 97 & 1.662e-4 & 1.841e-4 & 0.2866 & 0.01361 & 11.31\\
0.45 & 200 & 104 & 1.465e-5 & 1.647e-5 & 0.08888 & 2.205e-3 & 11.98\\
0.5 & 200 & 110 & 8.482e-7 & 9.506e-7 & 0.01650 & 1.752e-4 & 10.63\\
0.55 & 200 & 117 & 4.424e-8 & 6.066e-8 & 1.234e-3 & 1.213e-5 & 12.76\\
0.6 & 200 & 124 & 1.321e-9 & 1.813e-9 & 8.661e-5 & 4.921e-7 & 8.114\\
0.64 & 200 & 129 & 7.330e-11 & 1.036e-10 & 5.914e-6 & 3.008e-8 & 16.59\\
0.68 & 200 & 135 & 1.984e-12 & 2.832e-12 & 3.870e-7 & 1.180e-9 & 11.39\\
0.72 & 200 & 140 & 5.422e-14 & 1.207e-13 & 1.013e-8 & 3.167e-11 & 11.18\\
0.76 & 200 & 146 & 7.533e-16 & 1.802e-15 & 3.574e-10 & 7.456e-13 & 6.181\\
0.8 & 200 & 151 & 9.787e-18 & 3.681e-17 & 4.895e-12 & 1.086e-14 & 6.031\\
0.84 & 830 & 516 & 1.134e-19 & 1.757e-19 & 8.754e-14 & 1.322e-16 & 3.593\\
0.87 & 830 & 533 & 2.636e-21 & 4.329e-21 & 2.568e-15 & 3.892e-18 & 22.15\\
0.9 & 830 & 550 & 3.401e-23 & 5.545e-23 & 5.866e-17 & 5.819e-20 & 13.44\\
0.93 & 830 & 567 & 4.303e-25 & 7.594e-25 & 8.697e-19 & 9.536e-22 & 17.20\\
0.96 & 830 & 584 & 4.510e-27 & 7.555e-27 & 9.996e-21 & 8.439e-24 & 11.11\\
0.97 & 830 & 590 & 8.936e-28 & 1.613e-27 & 2.109e-21 & 1.929e-24 & 255.3\\
0.98 & 3610 & 2447 & 1.682e-28 & 2.165e-28 & 5.262e-22 & 4.471e-25 & 277.2\\
0.99 & 3610 & 2471 & 2.782e-29 & 3.552e-29 & 1.101e-22 & 7.684e-26 & 354.9\\
1.0 & 3610 & 2496 & 4.531e-30 & 6.601e-30 & 1.649e-23 & 1.141e-26 & 321.5\\
1.03 & 2450 & 1795 & 2.096e-32 & 4.110e-32 & 1.055e-25 & 6.604e-29 & 10.00\\
1.06 & 4000 & 2962 & 7.275e-35 & 1.249e-34 & 4.435e-28 & 2.691e-31 & 6.548\\
1.09 & 4000 & 3043 & 1.306e-37 & 2.162e-37 & 1.034e-30 & 4.711e-34 & 3.774\\
1.12 & 5550 & 4305 & 1.903e-40 & 3.001e-40 & 2.049e-33 & 8.502e-37 & 3.933\\
1.15 & 5550 & 4418 & 1.752e-43 & 2.799e-43 & 2.501e-36 & 8.572e-40 & 2.857\\
1.18 & 5550 & 4531 & 9.792e-47 & 1.636e-46 & 2.088e-39 & 6.011e-43 & 2.148\\
1.19 & 5550 & 4569 & 7.823e-48 & 1.312e-47 & 1.439e-40 & 4.329e-44 & 264.7\\
\bottomrule
\end{tabular}
\end{table}

At $L=0.8$ our certified values are consistent with Zhu's [Thm.~6.2]: $8.9\e{-18}\le\lambda_1^\ev\le2.523\e{-16}$, $\lambda_2^\ev\ge2.085\e{-12}$, $8.206\e{-15}\le\lambda_1^\od\le2.347\e{-14}$. Zhu's certified upper bound $\lambda_1^\ev(0.8)\le2.27\e{-17}$ [Table~3] is sharper than our $3.681\e{-17}$.

\subsection{Piece III}\label{sec:results3}
\emph{Test at $L=1.19$.} With Elston's published parameters at $L=1.19$ ($\delta=0.3$, $\beta_K=64$, $\Tsh=700$, $N=680$, $56$ Gauss points on panels of width $1/4$) the certifier, run by us, gives
\[
\lambda_1^\ev(1.19)\ge6.8131164696015\e{-48},\qquad \lambda_1^\od(1.19)\ge3.91966861609542239\e{-44},
\]
reproducing Elston's published $6.8131164696007\e{-48}$ (even, $12$ digits) and his odd bound to $17$ digits; the shift $\mu_s$ depends on the thread count through the numerical eigenvalue estimate, and both values are valid lower bounds. With \texttt{--deflate},
\[
\lambda_2^\ev(1.19)\ge1.2336395547\e{-40}
\]
(non-rigorous Galerkin value $1.23364079\e{-40}$). With $U(1.19)\le1.312\e{-47}$ from Table~\ref{tab:piece2}, $\min(\lambda_2^\ev,\lambda_1^\od)(1.19)/U(1.19)\ge2987$: simple-even at $L=1.19$ is confirmed by a second, independent code path (Elston's reduction instead of Zhu's pointwise cutoff, a different Gram assembly and a different Cholesky certificate).

\begin{table}[ht]
\centering\footnotesize
\caption{Piece III windows and certifier parameters, from the field \texttt{args} of each certificate file. Common to all runs: panel width $1/4$, Schur test on $4000$ cells with $1500$ weight iterations. ``GL'' is the number of Gauss--Legendre points per panel, ``bits'' the working precision, ``time'' the wall-clock seconds (field \texttt{seconds}) of the $\lambda_2^\ev$ (deflated) and $\lambda_1^\od$ runs, with $15$--$52$ worker processes. At $1.2785$ the two sectors used different parameters. New prime powers relative to the previous row are listed. Runs not listed are in progress.}\label{tab:grid3}
\begin{tabular}{lllrrrrrl}
\toprule
$L_i$ & $\delta$ & $\beta_K$ & $\Tsh$ & $N$ & GL & bits & time (s), $\lambda_2^\ev$ / $\lambda_1^\od$ & new $n$\\
\midrule
1.19 & 0.3 & 64 & 700 & 680 & 56 & 400 & 600--720 each & (test; $n\le 9$)\\
1.2205 & 0.3 & 64 & 850 & 849 & 56 & 400 & 1009 / 714 & 11\\
1.25 & 0.3 & 64 & 990 & 1010 & 56 & 400 & 1481 / 1034 & \\
1.2785, even & 0.3 & 64 & 1150 & 1196 & 57 & 400 & 2418 / --- & \\
1.2785, odd & 0.3 & 74 & 1130 & 1133 & 63 & 400 & --- / 1000 & \\
1.3055 & 0.3 & 78 & 1290 & 1315 & 65 & 400 & 3301 / 2466 & 13\\
1.3315 & 0.2 & 84 & 1570 & 1585 & 70 & 400 & 5673 / 4384 & \\
1.3565 & 0.2 & 89 & 1730 & 1757 & 74 & 600 & 11116 / 8545 & \\
1.3805 & 0.2 & 93 & 1890 & 1943 & 78 & 600 & 11639 / 7998 & \\
1.4035 & 0.2 & 98 & 2070 & 2143 & 81 & 600 & 14363 / 9784 & 16\\
\midrule
1.4255 & 0.2 & 102 & 2240 & 2356 & 85 & 600 & running / 13217 & 17\\
1.447 & 0.2 & 107 & 2440 & 2587 & 89 & 600 & 24322 / 20970 & \\
\bottomrule
\end{tabular}
\end{table}

\begin{table}[ht]
\centering\footnotesize
\caption{Piece III certificate. $U(L_i)\ge\lambda_1^\ev(L_i)$ is the $K=640$ Galerkin bound; $\ell^\od$ and $\ell_2$ are certified lower bounds for $\lambda_1^\od$ and $\lambda_2^\ev$; the step ratio in row $L_{i+1}$ is $\min(\ell_2,\ell^\od)(L_{i+1})/U(L_i)$. Lower bounds are rounded down and upper bounds up, to four significant digits, and ratios are floored. The lower bounds in row $1.19$ are the test of \S\ref{sec:results3} and enter no step. Below the rule: additional certified data points beyond $\Lmax$, not part of Theorem~\ref{thm:main} (the $\lambda_2^\ev$ run at $1.4255$ is in progress).}\label{tab:piece3}
\begin{tabular}{llllll}
\toprule
$L_i$ & $U(L_i)$ & $\lambda_1^\od\ge$ & $\lambda_2^\ev\ge$ & step ratio $\ge$\\
\midrule
1.19 & 8.712e-48 & 3.919e-44 & 1.233e-40 & (Piece II: 264.7)\\
1.2205 & 2.336e-51 & 1.224e-47 & 4.200e-44 & 1.405\\
1.25 & 4.458e-55 & 2.642e-51 & 1.162e-47 & 1.131\\
1.2785 & 8.317e-59 & 5.889e-55 & 2.648e-51 & 1.321\\
1.3055 & 1.209e-62 & 9.466e-59 & 4.806e-55 & 1.138\\
1.3315 & 1.693e-66 & 1.706e-62 & 9.901e-59 & 1.412\\
1.3565 & 1.962e-70 & 2.187e-66 & 1.435e-62 & 1.292\\
1.3805 & 2.251e-74 & 3.032e-70 & 2.247e-66 & 1.546\\
1.4035 & 2.493e-78 & 3.492e-74 & 2.991e-70 & 1.552\\
\midrule
1.4255 & 2.371e-82 & 3.564e-78 & (running) & ---\\
1.447 & 2.107e-86 & 3.269e-82 & 3.420e-78 & 1.379\\
\bottomrule
\end{tabular}
\end{table}

% =====================================================================
\section{Rigorous upper bounds for larger windows}\label{sec:upper}

This section gives \emph{upper bounds only}; it makes no lower-bound, positivity or simplicity statement.

\emph{Galerkin matrix.} Let $e_k(y)=\sqrt{(2k+1)/(2L)}P_k(y/L)$ with $k=2m+p$, $m<K$, $p\in\{0,1\}$ the parity. With $r=x/L\in[0,2]$, the symmetrised autocorrelation of $e_a,e_b$ is an $L$-independent polynomial in $r$: up to the normalisation $\sqrt{(2a+1)(2b+1)}/2$ it equals $C_{ab}(r)=\int_{-1}^{1-r}P_a(t)P_b(t+r)\,dt$ (for $a+b$ even $C_{ab}=C_{ba}$). By~\eqref{eq:P} the Galerkin matrix is $Q_K=\bigl(\int_0^2C_{ab}\,d\nu_L\bigr)_{a,b}$ plus a multiple of the identity, where $\nu_L$ is the measure with smooth density (pole and archimedean weights) minus the point masses $2\Lambda(n)n^{-1/2}$ at $r=\log n/L$. The entries follow from the Legendre moments $\int P_j(r-1)\,d\nu_L$ by the convolution recurrence~\cite{HaleTownsend}
\[
C_{a+1,b}=C_{a-1,b}+\frac{2a+1}{2b+1}\bigl(C_{a,b-1}-C_{a,b+1}\bigr),
\]
with exact seeds $C_{0b},C_{1b}$ expressed through $\int_u^1P_k$. The recurrence is run only on $b\ge a$, where it is forward stable; the cost is $O(K^2)$ instead of $O(K^4)$ for direct quadrature.

\begin{proposition}[Rigorous moments]\label{prop:upper}
The smooth moments are computed with $M_q$-point Gauss--Legendre in $r$, and the truncation error of each is bounded by Lemma~\ref{lem:G}: $|P_j|\le\rho^j$ on $E_\rho$, the maxima of the weights $4L\cosh(Lr/2)$, $2Lr\,e^{-Lr/2}/(1-e^{-2Lr})$ and of the diagonal weight are enclosed by covering the ellipse with $2000$ complex balls, and the poles $r=i\pi m/L$, $m\ne0$, are checked to lie outside $E_\rho$. The Legendre values at the nodes are computed by the Chebyshev recurrence on exact midpoints with an a priori rounding bound and a Markov bound for the node radii, followed by the nonnegative Chebyshev-to-Legendre map. Hence every entry of $Q_K$ is enclosed in a ball, and for every exact vector $v$ the upper endpoint of $v^TQ_Kv/v^Tv$ is an upper bound for $\lambda_1^\sigma(L)$ in the sector of~$v$.
\end{proposition}
The last statement holds because $v$ defines a polynomial trial function in the form domain (\S\ref{sec:piece2cert}). The vector $v$ comes from inverse iteration on midpoints and carries no rigour; the reported value is outward-rounded.

\begin{table}[ht]
\centering\footnotesize
\caption{Rigorous upper bounds for the lowest eigenvalue in each sector. $K$ = number of Legendre modes per sector; ``---'' = no rigorous run. The Piece III windows are in Table~\ref{tab:piece3} ($K=640$, $3000$ bits).}\label{tab:upper}
\begin{tabular}{llll}
\toprule
$L$ & $\lambda_1^\ev(L)\le$ & $\lambda_1^\od(L)\le$ & $K$, bits (even; odd)\\
\midrule
1.5 & 3.85105826960019e-97 & 8.21780562022764e-93 & $384$, $2400$; $512$, $4000$\\
1.75 & --- & 2.06808416452444e-162 & ---; $512$, $4000$\\
2.0 & 1.84681871820091e-283 & 3.23821041502563e-278 & $1024$, $6000$; $1024$, $6000$\\
2.25 & 9.03204367909380e-476 & 4.64181809939520e-470 & $2048$, $8000$; $1792$, $8000$\\
\bottomrule
\end{tabular}
\end{table}

At $L=2$ our even bound is below Zhu's certified $3.19\e{-283}$ [Table~3]; certified upper bounds for $L\le2$ are not new. The window $L=2.25$ lies outside the range of~\cite{Zhu}. There ($\mu=e^{2L}\approx90.0$) the explicit analytic bounds we know of are larger: about $e^{-1052.9}$ from~\cite[Thm.~K(a)]{MoriV} (in his normalisation) and about $e^{-504.0}$ from~\cite[Thm.~2.1]{rh25}, against $e^{-1093.8}$ for our even bound; \cite[Thm.~1]{rh81} assumes RH and has a non-explicit $O(L^2)$ term, so it gives no value. The values of Table~\ref{tab:upper} at $L=2.25$ are therefore the sharpest explicit upper bounds there that we know of; the decay scale $-\ln\lambda\approx4\pi e^{2L}$ is not ours~\cite{Connes26,MoriV,rh80,rh81}. Non-rigorous Galerkin values continue to decrease slowly with $K$ (at $L=2.25$: $9.894\e{-476}$, $9.337\e{-476}$, $9.032\e{-476}$ for $K=1536,1792,2048$); we do not have a bound on the omitted-mode tail, so these upper bounds may be far from sharp in relative terms.

\begin{remark}[Quadrature exponent]\label{rem:rho}
The $L=2$ even run was made before the exponent of Lemma~\ref{lem:G} was corrected from $\rho^{-2n}$ to $\rho^{-2n+2}$ in this code path ($\rho\le3$ there). The quadrature contributions to the moment radii are below $10^{-1999}$ and the reported radius of the Rayleigh quotient is $4.5\e{-1656}$; scaling the quadrature radii by $\rho^2\le9$ cannot reach the slack of the $15$-digit outward rounding (order $10^{-297}$). A re-run with the corrected exponent has not been made. The bounds $U$ of Table~\ref{tab:piece3} were computed on 30 October 2026 with \file{fr_sweep.py} as shipped in \file{code/upper/}: the exported value is the upper end of the Rayleigh-quotient ball raised by a relative $10^{-13}$ before rounding to $15$ digits, while every ball radius is below $10^{-600}$ times the value, so a factor $\rho^2\le9$ on the quadrature part of the radius cannot reach that margin, whichever exponent was in use. The working directory was not under version control, so the revision of \file{fast_recurrence.py} used for each run is not recorded. Re-running $U(1.19)$ with the shipped code on 31 October 2026 reproduced every exported digit and the radius (\file{data/piece3/rerun_U1.19.jsonl}). For the other rows of Table~\ref{tab:upper} (28--30 October; by-products that enter no proof) we have not confirmed which exponent was used.
\end{remark}

% =====================================================================
\section{A numerical observation: prime-power singular points}\label{sec:kinks}

\emph{Everything in this section is numerical and non-rigorous.} It concerns finite-dimensional approximants of the even minimiser; whether the limiting eigenfunction has the same features is not established. That near-null eigenvectors have weak singularities at $\log n$ and $L-\log n$, and that meshes should be cut there, is due to Eberwein~\cite[\S13]{Eberwein}, whose spectral-element code already does so; the mechanism is the one familiar from differential-difference equations. Only the measurements below are ours. Note also that second-generation points such as $L-\log6$ exist, so the singular set is not limited to $|L-\log n|$ with $n$ a prime power.

At $L=1.19$ the Legendre coefficients of the lowest Galerkin eigenvector ($K=384$) fall exponentially to about $10^{-25}$ by degree $160$ and then only like (degree)$^{-2.1}$. A periodogram of the high-degree coefficients locates singular points at $|y|\approx0.088$--$0.105$, $0.183$--$0.198$, $0.406$, $0.510$, $0.747$, $0.884$, $1.003$, which match $|L-\log n|=0.091$, $0.196$, $0.419$, $0.497$, $0.756$, $0.889$, $1.007$ for $n=3,4,5,2,7,8,9$; there is no signal at $|L-\log6|=0.602$. A heuristic reason: in the eigen-equation the shifted copies $f(y\pm\log n)$ (with $f=0$ outside $[-L,L]$) switch on or off where $y\pm\log n$ crosses $\mp L$, and only prime powers carry weight ($\Lambda(6)=0$).

To quantify this we use a split-window basis (piecewise Legendre, discontinuous at $\pm(L-\log n)$; its $x$-quadrature error is checked by node doubling but not bounded) and record the one-sided jumps of the computed eigenvector at the breakpoints, normalised by $\Lambda(n)n^{-1/2}f(L)$. Table~\ref{tab:kinks} gives the absolute values at $L=1.19$ for four resolutions. The signs are negative for $n=2,3$ (where $L-\log n>0$) and positive otherwise, as the heuristic predicts.

\begin{table}[ht]
\centering\footnotesize
\caption{$|\text{jump}|/(\Lambda(n)n^{-1/2}f(L))$ at $y=L-\log n$ (or $\log n-L$), $L=1.19$, even sector, split-window basis with base allocation $24$; ``CV'' is the coefficient of variation over the seven prime powers.}\label{tab:kinks}
\begin{tabular}{lrrrrrrrrr}
\toprule
$N_{\mathrm{tot}}$ & $n=3$ & 4 & 5 & 2 & 7 & 8 & 9 & mean & CV\\
\midrule
600 & 0.119 & 0.134 & 0.112 & 0.117 & 0.112 & 0.141 & 0.120 & 0.122 & 0.085\\
800 & 0.112 & 0.125 & 0.104 & 0.117 & 0.107 & 0.120 & 0.110 & 0.113 & 0.061\\
1200 & 0.107 & 0.109 & 0.103 & 0.107 & 0.104 & 0.119 & 0.118 & 0.110 & 0.054\\
1600 & 0.100 & 0.108 & 0.099 & 0.106 & 0.097 & 0.112 & 0.107 & 0.104 & 0.051\\
\bottomrule
\end{tabular}
\end{table}

At $L=1.19$ the normalised jumps of all seven prime powers are thus close to a common value (CV about $5\%$ at the two highest resolutions), although raw weights $\Lambda(n)n^{-1/2}$ range from $0.25$ to $0.74$; the common value drifts down with resolution (roughly like $1/(\ln(1/h)+b)$ in the mesh size $h$), so it is not a constant: the computed jumps are finite-resolution images of a logarithmically softened singularity of a continuous function, and they tend to $0$. At $L=1.5$ the value-jump ratios are again close to each other, but the split-window mesh there omits some cuts (including second-generation points close to the $n=8,16$ cuts), and we make no per-breakpoint claim. Prior work: the location of the singular points and the cut mesh are Eberwein's~\cite[\S13]{Eberwein}; jumps of a trial vector at $L-\log n$ appear in Mori's work~\cite{Mori}, and the $\pm\log n$ shift structure in~\cite[eq.~(2.11)]{Suzuki}.

% =====================================================================
\section{Reproducibility}\label{sec:repro}

\subsection{Software and hardware}
Python 3.12 (3.12.13 for Pieces I--II) with python-flint 0.9.0, which bundles FLINT/Arb 3.6.0~\cite{FLINT,Johansson}; the C tools \texttt{vchol} (verified Cholesky) and \texttt{veig} (estimates only) of Piece~II are built with gcc 15.2.0 against FLINT 3.4. Pieces I--II ran on a 20-thread laptop and a cluster of three 64-core nodes (FLINT versions on the cluster not recorded; a SHA-256 check of the code files on every host was in force for the certificate of record). Piece~III and \S\ref{sec:upper} ran on three 64-vCPU virtual machines. Elston's certifier is used at upstream commit \file{3cc4cb7be6a659f1463ae4e2e99ad7ae5e6646f5} with the \texttt{--deflate} option added (\file{code/gap/elston/}).

\emph{Run times.} Piece~II certificate of record: clean rerun for $L\le1$ about $5.8$\,h in three stages; extension $1.03$--$1.09$: $2859$\,s, $11022$\,s, $11760$\,s; continuation $1.12$--$1.19$ about $28.5$\,h in total (head assembly about $6.5$--$7$\,h per window on the cluster). Piece~III test at $L=1.19$: $10$--$12$\,min per run with $18$ workers. Piece~III certificates: per run as in Table~\ref{tab:grid3}; the $16$ runs behind Theorem~\ref{thm:main} took about $814$ worker-process hours (seconds times workers), all $19$ shipped runs about $1177$. Upper bounds $U$ ($K=640$, $3000$ bits): $626$--$822$\,s per window. Rigorous upper bounds of Table~\ref{tab:upper}: $L=2$, $K=1024$: $3166$\,s (even), $10971$\,s (odd); $L=2.25$: $63923$\,s (even), $77675$\,s (odd); $L=1.5$, $1.75$ odd: $1278$\,s, $1293$\,s; $L=1.5$ even ($K=384$): $421$\,s.

\subsection{Commands}
\begin{itemize}
\item Piece I and Piece II gluing: \texttt{python code/theorem.py --check} (seconds; ends with ``THEOREM CERTIFIED for all L in (0, 1.19]''). Windows: \texttt{python code/theorem.py} with the environment settings recorded in the release.
\item Piece III gluing: \texttt{python code/gap/check\_steps.py} (seconds; reads \file{data/piece3/}; ends with ``CERTIFIED: \dots simple-even for all L in (0, 1.4035]'', compare \file{data/piece3/check.log}).
\item Piece III lower bounds, in \file{code/gap/elston/}, with the window options \texttt{<w>} from the field \texttt{args} of the certificate file (e.g.\ \texttt{--a 1.2205 --delta 0.3 --beta-k 64 --tsharp 850 --nmodes 849 --gl 56 --prec 400}):
{\footnotesize
\begin{verbatim}
C="--panels-per-unit 4 --schur-k 4000 --schur-iters 1500"
python grid_certificate_rigorous.py <w> $C --sector odd --workers n --json <out>
python grid_certificate_rigorous.py <w> $C --sector even --deflate --workers n \
       --json <out>
\end{verbatim}}
\item Upper bounds, in \file{code/upper/} (exports the outward-rounded field \texttt{lam0\_upper}):
{\footnotesize
\begin{verbatim}
THREADS=n python fr_sweep.py <out.jsonl> "640" "1.19,1.2205,..." 3000
\end{verbatim}}
\end{itemize}

\begin{table}[ht]
\centering\footnotesize
\caption{Provenance of the numbers in this paper. Paths under \texttt{highprec/} are in the author's working tree and are not part of the repository; all files the theorem depends on are.}\label{tab:files}
\begin{tabular}{>{\raggedright\arraybackslash}p{0.36\textwidth}>{\raggedright\arraybackslash}p{0.56\textwidth}}
\toprule
Item & File(s)\\
\midrule
Table~\ref{tab:piece2}, Piece I constants & \file{data/theorem_bounds.json}; \file{code/theorem.py}; logs \file{data/logs/check.log}, \file{run_rerun.log}, \file{run_extend.log} (release of~\cite{Botha})\\
Remark~\ref{rem:erratum} & \file{data/history/theorem_bounds_corrected.json}, \file{data/logs/run_gl_correction.log}\\
Piece III parameters (Table~\ref{tab:grid3}) & field \texttt{args} of each file in \file{data/piece3/}; design rule \file{code/gap/params.py}\\
Piece III certifier & \file{code/gap/elston/grid_certificate_rigorous.py}, \file{UPSTREAM_COMMIT}, \file{LICENSE}\\
$L=1.19$ test (\S\ref{sec:results3}) & run logs on the compute host, not shipped (the test enters no step of the proof)\\
Table~\ref{tab:piece3} lower bounds & \file{data/piece3/g_<L>_l2.json} ($\lambda_2^\ev$), \file{g_<L>_odd.json} ($\lambda_1^\od$); field \texttt{certified\_lower\_bound}\\
Table~\ref{tab:piece3} upper bounds & \file{data/piece3/upper640.jsonl}, field \texttt{lam0\_upper}; code \file{code/upper/}\\
Piece III gluing & \file{code/gap/check_steps.py}, log \file{data/piece3/check.log}, checksums \file{data/piece3/SHA256SUMS}\\
Table~\ref{tab:upper}, $L=2$ even & \file{highprec/data/rig_L2.jsonl}, \file{rig_L2.log}\\
Table~\ref{tab:upper}, other rows & \file{highprec/data/rig_L225.jsonl} ($L=2.25$ even), \file{highprec/data/odd.jsonl} (odd sector), \file{highprec/gap/out_all/upper.jsonl} ($L=1.5$ even)\\
Non-rigorous Galerkin values (\S\ref{sec:upper}) & \file{highprec/data/fr_L225*.jsonl}, \file{fr_sweep*.jsonl}\\
\S\ref{sec:kinks} & \file{highprec/data/tail_probe.json}, \file{kink_119_*.json}, \file{kink_15_*.json}\\
\bottomrule
\end{tabular}
\end{table}

\subsection{Verification and audits}
All reviews so far were carried out by AI agents directed by the author; no independent human referee has checked the proofs or the code.
\begin{itemize}
\item \emph{Pieces I--II.} Three independent review rounds of the earlier note produced, among others: the Gauss--Legendre exponent correction (Remark~\ref{rem:erratum}); reconciliation of the Piece I constants across precisions; outward rounding of all quoted bounds; a downward-induction proof of the Bessel ratio enclosure; the definition of the continuation of $\re\psi$ off the real axis; documentation of Python \texttt{min} on balls; a mollification argument for~\eqref{eq:P}; a self-contained discreteness proof; a SHA-256 check of remote code; and corrected credit to Zhu for the reduction. Every finding was checked against code, logs and data before it was applied. A clean end-to-end rerun with a single code revision (empty cache, tools rebuilt, remote code hashed) then reproduced all lower bounds for $L\le1$ to about three significant digits. The sharp-cutoff path and the windows $1.03$--$1.19$ have not been through a review round.
\item \emph{Upper-bound code.} An external review of the high-precision notes (six findings, all accepted) led to per-code-path rigour labels, outward rounding of exported upper endpoints, and withdrawal of an unsupported tail claim; an earlier review found and fixed an inaccurate ball inverse, a missing residual check and a missing quadrature bound.
\item \emph{Piece III.} Elston's certifier was audited independently at $L=1.19$ (two issues fixed; documented in~\cite{Elston}); the post-audit code at larger $L$ and our \texttt{--deflate} addition have not been audited. The step test (\file{code/gap/check_steps.py}, about $50$ lines) compares explicit Arb endpoints. The reproduction of Elston's $L=1.19$ bounds to $12$ and $17$ digits is a consistency check of the installation, not an audit.
\end{itemize}
\emph{Invitation to review.} This work has not been peer-reviewed. Code, certificate files and this source are public at \url{https://github.com/AlexGenusCore/weil-ground-state-simple-even}; the file \file{REVIEWING.md} there lists the least-checked parts and concrete checks ordered by effort, from re-running the gluing script on the shipped certificates (minutes) to re-running a single window (hours). Independent review by anyone is welcome; confirmed errors will be corrected publicly and reviewers credited.

% =====================================================================
\section{Limitations}\label{sec:limitations}
\begin{enumerate}
\item \emph{Range.} The theorem is certified on $(0,\Lmax]$. The extension to $1.5$ is in progress (Remark~\ref{rem:status}); the passing step $[1.4255,1.447]$ is a data point only.
\item \emph{Piece III is unaudited beyond $L=1.19$.} Elston's certifier was audited only at $L=1.19$; the runs at larger $L$, the parameter changes ($\beta_K$, $\Tsh$, $N$, Gauss points) and our \texttt{--deflate} addition have not been reviewed. The step ratios of Piece~III are small ($1.131$--$1.552$), so the theorem on $(1.19,\Lmax]$ depends on every bound being correct to within about $13\%$.
\item \emph{Computer assistance.} Correctness depends on Arb/FLINT, on our code for Pieces I--II and \S\ref{sec:upper}, on Elston's code for Piece~III, and on the lemmas above.
\item \emph{Normalisation.} Our $Q_L$ follows Zhu and was validated against the zero side of the explicit formula. It has not been checked line by line against $QW_\lambda$ of~\cite{CvS,CCM25}. Remark~\ref{rem:affine} covers positive affine changes, not a discrepancy in the relative weights of the pole, archimedean and prime terms.
\item \emph{Self-adjoint realisation.} We use the closure of $Q_L$ on $C_c^\infty(-L,L)$; trial polynomials not vanishing at $\pm L$ are in its form domain by the argument in \S\ref{sec:piece2cert}. A different realisation would need re-examination.
\item \emph{Unexplained precision dependence} of the upper end of $\int|h|$ in Piece~I (largest value used).
\item \emph{Comparisons on balls.} Python \texttt{min}/\texttt{max} on Arb balls occur in the gluing and in some intermediate steps of both certifiers; every current use is separated by many orders of magnitude more than the radii.
\item \emph{Sharp-cutoff and triangular assembly code} for $1.03\le L\le1.19$ in Piece~II has not been reviewed; $L=1.19$ is, however, re-confirmed by the Piece~III code path.
\item \emph{Upper bounds of \S\ref{sec:upper}} are upper bounds only; their distance from the true eigenvalue is not controlled (Remark~\ref{rem:rho} for the $L=2$ quadrature exponent).
\item \emph{Section~\ref{sec:kinks}} is numerical; the split-window $x$-quadrature is not bounded.
\item \emph{Scope.} Nothing is claimed for $L>\Lmax$, about the convergence step~(ii) of~\cite{CCM25}, or about RH.
\end{enumerate}

\section*{Acknowledgements}
The Piece~III lower bounds use Zac Elston's certifier \emph{weil-windows}~\cite{Elston} (MIT licence), which we ran with one added option; the reduction it implements and the published $L=1.19$ certificates are his. The window certificates of Pieces I--II rest on Xuefeng Zhu's frequency-cutoff reduction and block estimate~\cite{Zhu}. The survey of lower-bound methods behind Piece~III drew on the work of Mori~\cite{Mori} and of Liu~\cite{Liu}, whose position-side use of the Legendre convolution recurrence preceded ours.

\section*{AI-assistance statement}
Computations, code and drafting were carried out with substantial assistance from an AI system (Anthropic's Claude). All results were checked by independent recomputation, and the author takes responsibility for the content.

\bibliographystyle{amsplain}
\bibliography{refs}

\end{document}
